
Comparar a viabilidade técnica, a acurácia e o consumo energético entre o processamento de \gls{ocr} realizado localmente no \gls{esp32cam} (abordagem Edge Computing) e a inferência \gls{ocr} realizada remotamente em servidor na nuvem (Cloud), utilizando a rede \gls{lorawan} como meio de transmissão.

\subsection{Objetivos secundários}

\begin{enumerate}
    \item Implementar os dois fluxos de processamento que compõem o experimento comparativo: o fluxo Edge, em que o \gls{ocr} ocorre localmente no \gls{esp32cam}; e o fluxo Cloud, em que a imagem é transmitida para processamento remoto. % 1. construção do experimento
    \item Configurar a rede \gls{lorawan} conectando os nós da rede com o objetivo de reduzir ao máximo o consumo energético; % 2. Configuração da infrastrutura de rede
    \item Avaliar e comparar os dois fluxos segundo as métricas: precisão de leitura (\%), latência de processamento (ms), consumo energético por leitura (mAh) e volume de dados transmitidos (bytes/leitura), identificando o \textit{trade-off} entre processamento local e transmissão. % 3. Responde a pergunta central do TCC
\end{enumerate}